\documentclass[11pt]{article}
\usepackage[utf8]{inputenc}
\usepackage[T1]{fontenc}
\usepackage{booktabs}
\usepackage{amsmath}
\usepackage{geometry}
\geometry{margin=2.5cm}

\title{Causal Fairness Analysis Report}
\author{Generated by FairMind and an LLM}
\date{2026-08-22}

\begin{document}
\maketitle

\section{Analysis setup}

\begin{tabular}{ll}
\toprule
Dataset & german\_credit \\
Protected attribute ($X$) & S2\_Sex \\
Comparison & Female $\to$ Male \\
Outcome ($Y$) & T\_Creditability (1) \\
Mediator ($W$) & Length of current employment \\
Confounder ($Z$) & Occupation \\
\bottomrule
\end{tabular}

\section{Causal effects (FairMind, exact values)}

The five effects below are computed by FairMind through exact inference on the
fitted Bayesian network. These values are final and must not be recomputed or
altered.

\begin{tabular}{lr}
\toprule
Effect & Value \\
\midrule
Total Variation (TV) & 0.077323 \\
Total Effect (TE) & 0.078572 \\
Spurious Effect (SE) & -0.001250 \\
Direct Effect (DE) & 0.050708 \\
Indirect Effect (IE, decomposition form: TE = DE + IE) & 0.027864 \\
\bottomrule
\end{tabular}

\section{Qualitative interpretation}

\subsection{Total effect and spurious component (TV, TE, SE)}

The total disparity between the two groups is 0.077323, indicating a significant difference in creditability outcomes. After removing the confounding variable, the genuine causal effect (TE) is 0.078572, which is slightly larger than the total variation, suggesting that the confounder slightly reduced the observed disparity. The spurious effect (SE) is -0.001250, indicating that a small portion of the observed disparity was driven by the confounder rather than by the protected attribute itself.

\subsection{Direct effect (DE)}

The direct effect (DE) is 0.050708, indicating a clear direct discrimination against the protected group (females). This direct effect is substantial, as it is more than 5\% of the total effect, suggesting that the protected group faces a significant disadvantage that is not mediated through the employment length.

\subsection{Indirect effect (IE)}

The indirect effect (IE) is 0.027864, which amplifies the direct effect since it has the same sign. This means that the mediator (employment length) adds to the overall discrimination, increasing the total effect from the direct effect alone. The indirect effect is positive, indicating that the mediator channel does not mitigate but rather exacerbates the direct discrimination.

\section{Recap Questions}

\begin{enumerate}
\item Is there direct discrimination against the protected group? \textbf{LLM:} YES \quad \textbf{Expected:} YES \quad (correct)
\item Does the indirect effect through the mediator mitigate the total effect? \textbf{LLM:} NO \quad \textbf{Expected:} NO \quad (correct)
\item Does the observed total effect exceed the practical relevance threshold? \textbf{LLM:} YES \quad \textbf{Expected:} YES \quad (correct)
\item Does the spurious component (SE) contribute substantially to the observed difference? \textbf{LLM:} NO \quad \textbf{Expected:} NO \quad (correct)
\item Is the direct effect the dominant component compared to the indirect effect? \textbf{LLM:} YES \quad \textbf{Expected:} YES \quad (correct)
\end{enumerate}

\vspace{1em}
\noindent\footnotesize
The recap answers follow deterministic threshold rules applied to the values above: direct discrimination when $|\mathrm{DE}| > 0.05$; a mediator channel counted as negligible when $|\mathrm{IE}| < 0.005$; practical relevance when $|\mathrm{TV}| > 0.05$; a substantial spurious component when $|\mathrm{SE}| / |\mathrm{TV}| > 0.1$.


\vspace{1em}
\noindent\textbf{Answer consistency:} 5/5 (100.0\%). The expected answers are derived deterministically from the FairMind values alone, through threshold rules, and were added to the document after it was generated: the model never saw them.

\end{document}